Monte Carlo output is only evidence about a phase transition if two things hold: the Markov chain samples the distribution we intend, and the observable we read off identifies the transition we intend. We examine both for the two-dimensional Ising lattice gas at fixed composition under pair-exchange (Kawasaki-type) Metropolis dynamics, used as a lattice model of residential segregation. First, we derive the exact exchange energy, which differs from the sum of two single-flip energies by $4J$ when the exchanged sites are neighbours, and prove that the resulting kernel is reversible, irreducible and aperiodic with the canonical law as its unique stationary distribution. Dropping the adjacency term breaks detailed balance by exactly $e^{\pm4\beta J}$; with uniform random-pair proposals the one-step perturbation is at most $4N^{-2}(1-e^{-4\beta J})$. A small one-step perturbation does not by itself bound the stationary one: on tori small enough to solve exactly the stationary law moves by $\EtwoPairMin$--$\EtwoNNMax$ in total variation, and at production sizes the mean-energy shift is a few per cent at $N=8,16$ and not detectable with our replication at $N\ge32$. Second, we compare the energy-variance peak, a common estimate of a transition or demixing temperature, with the exact fixed-composition coexistence curve obtained from Onsager and Yang. Over eight replicates, the single-run estimate has a standard deviation of $\EfiveSdMin$--$\EfiveSdMax$, larger than the whole composition dependence of the exact curve over $0.2\le f\le0.8$ ($\EbDevTwoTenths$, or $\EbRelTwoTenths\%$). At dilute compositions it lies $0.17$--$0.29$ below the exact curve. A long equilibrium reference shows that the standard 300-sweep sampling window is shorter than the energy autocorrelation time in the worst cases, but this does not account for the dilute offset, which persists in the equilibrium reference and is consistent with finite-size droplet formation. Tables, figures and the numbers in them are generated from stored raw output by the accompanying code.
